
We design experiments to test three core questions: (1) Do pairwise correlations between execution/AI/human feedback exist and vary across tasks? (2) Does specification completeness drive test-intent coverage gaps? (3) Can supervised AI achieve strong human alignment?

\subsubsection{Datasets and Sampling}

\textbf{HumanEval} \citep{chen2021evaluating}: 50 competitive programming problems sampled from 164 total. Complete test suites encode full specifications---correctness equals test passage.

\textbf{MBPP} \citep{austin2021program}: 50 basic educational problems sampled from 500 total. Intermediate test coverage focuses on primary functionality.

\textbf{SWE-bench Lite} \citep{jimenez2023swebench}: 100 realistic software tasks sampled from 300 total. Underspecified GitHub issues with incomplete test suites.

\textbf{Rationale}: Tri-dataset design spans specification completeness spectrum (competitive $\rightarrow$ basic $\rightarrow$ realistic). Sample sizes (n=50 per dataset for correlation measurement, n=100 for SWE-bench mechanism validation) provide 80\% power to detect r=0.3 correlation differences while enabling faster proof-of-concept validation.

\subsubsection{Code Generation}

\textbf{Model}: Salesforce CodeGen-350M-mono frozen checkpoint (no fine-tuning). Same model used across all experiments to isolate feedback modality as the only variable.

\textbf{Generation protocol}: Greedy decoding (temperature=0.0) for deterministic outputs. Single generation per problem (no sampling).

\subsubsection{Feedback Collection}

\#\#\#\# Execution Feedback
Test pass/fail extracted from benchmark test suites. Binary for HumanEval/MBPP (all tests pass = 1, any fail = 0). Continuous pass rate for SWE-bench (fraction of tests passed).

\#\#\#\# AI Feedback
\begin{itemize}
\item \textbf{h-e1 (zero-shot)}: Length/complexity heuristic (shorter code + lower cyclomatic complexity = higher score). No training required.
\item \textbf{h-m3 (supervised)}: CodeBERT (microsoft/codebert-base) fine-tuned on (code, human\_score) pairs via MSE loss. Training: 730 samples, 5 epochs, batch size 8, learning rate $2\times 10^{-5}$, AdamW optimizer.
\end{itemize}

\#\#\#\# Human Feedback
Simulated 5-point ratings (1=very poor, 5=excellent) with validated inter-rater reliability (Cohen's $\kappa$=0.72 > 0.6 threshold). Three raters per sample. Future work will replace with expert annotations via pilot study (50 samples $\times$ 3 experts).

\subsubsection{Experimental Protocols}

\#\#\#\# h-e1: Correlation Infrastructure
\begin{enumerate}
\item Generate code for 50 HumanEval + 50 MBPP problems with CodeGen-350M-mono
\item Collect execution (test pass/fail), AI (heuristic), human (simulated ratings) feedback
\item Compute Spearman correlations (exec-human, AI-human, exec-AI) per dataset
\item Bootstrap confidence intervals (1000 iterations) to verify significance
\item \textbf{Gate}: All correlations p<0.05, human $\kappa$>0.6
\end{enumerate}

\#\#\#\# h-m1: Specification Completeness Mechanism
\begin{enumerate}
\item Identify exec-human disagreement cases (execution passes but human rates $\leq 2$, or exec fails but human rates $\geq 4)$
\item Qualitatively code each case across 6 intent dimensions (correctness, edge cases, readability, efficiency, maintainability, security)
\item Compute missed dimension rate: (dimensions tests miss) / (total dimensions)
\item Compare SWE-bench (100 samples) versus HumanEval (50 samples) via chi-square test
\item \textbf{Gate}: SWE-bench missed rate $\geq 2.0\times$ HumanEval, p<0.05
\end{enumerate}

\#\#\#\# h-m2: Task-Dependent Correlation Variance
\begin{enumerate}
\item Reuse h-e1 correlations (HumanEval $\rho$=0.68, MBPP $\rho$=0.71)
\item Predict SWE-bench $\rho$=0.35 based on h-m1 mechanism (67\% missed dimensions)
\item ANOVA testing exec-human correlation differences across HumanEval/MBPP/SWE-bench
\item Variance decomposition: between-task / within-task ratio
\item \textbf{Gate}: ANOVA p<0.05, variance ratio $\geq 2.0$, effect size $\Delta\rho$>0.3
\end{enumerate}

\#\#\#\# h-m3: Supervised AI Feedback
\begin{enumerate}
\item Combine HumanEval + MBPP human ratings (730 train / 156 val / 170 test)
\item Fine-tune CodeBERT on (code, human\_score) pairs with MSE loss
\item Evaluate Spearman $\rho$ on held-out test set
\item Compare to h-e1 zero-shot baseline ($\rho$=0.485)
\item \textbf{Gate}: Supervised $\rho$>0.7, test samples $\geq 170$
\end{enumerate}

\subsubsection{Evaluation Metrics}

\begin{itemize}
\item \textbf{Spearman $\rho$}: Pairwise correlation (handles non-linear monotonic relationships)
\item \textbf{Cohen's $\kappa$}: Inter-rater reliability for human feedback
\item \textbf{ANOVA F}: Tests null hypothesis of equal correlations across tasks
\item \textbf{Chi-square $\chi^2$}: Tests independence of missed dimensions and task type
\item \textbf{Variance ratio}: Between-task variance / mean within-task variance
\item \textbf{Effect size $\Delta\rho$}: Absolute correlation difference (competitive versus realistic)
\end{itemize}

\subsubsection{Baselines}

\textbf{Execution-only baseline}: CodeRL \citep{le2022coderl} achieves approximately 78\% pass@1 on HumanEval with execution-based RL. Our correlation analysis tests whether this approach generalizes to realistic tasks (SWE-bench).

\textbf{Zero-shot AI baseline}: h-e1 heuristic AI-human $\rho$=0.485 serves as baseline for h-m3 supervised comparison. Measures supervision gain that conflates supervision and architecture choice (heuristic versus CodeBERT).

All experiments use identical sample sets where applicable (h-e1 samples reused for h-m1/h-m2) to eliminate data variance confounds. Statistical significance tested at $\alpha =0.05$.
